\section*{2 Partial Derivatives}

\subsection*{2.1 Definition of Partial Derivatives}

\begin{conceptbox}
    A partial derivative measures the rate of change of a multivariable function with respect to one variable while holding all other variables constant:
    \begin{equation*}
        \frac{\partial f}{\partial x} = \lim_{h \to 0} \frac{f(x+h, y, z) - f(x, y, z)}{h}
    \end{equation*}
\end{conceptbox}

Partial derivatives extend the idea of ordinary derivatives to higher dimensions. Instead of a slope along a line, we now measure slopes along coordinate directions.

\begin{examplebox}
    Find $\frac{\partial f}{\partial x}$ for $f(x,y)=x^2y+3y^2$.
    \newline \textbf{Solution:}
    \begin{align*}
        \frac{\partial f}{\partial x} & = \frac{\partial}{\partial x}(x^2y+3y^2) \\
                                      & = 2xy + 0                                \\
                                      & = 2xy
    \end{align*}
\end{examplebox}

\begin{examplebox}
    Find $\frac{\partial f}{\partial y}$ for the same function.
    \newline \textbf{Solution:}
    \begin{align*}
        \frac{\partial f}{\partial y} & = \frac{\partial}{\partial y}(x^2y+3y^2) \\
                                      & = x^2 + 6y
    \end{align*}
\end{examplebox}

\begin{examplebox}
    Evaluate $\frac{\partial f}{\partial x}$ at $(1,2)$.
    \newline \textbf{Solution:}
    \begin{align*}
        \frac{\partial f}{\partial x} = 2xy = 2(1)(2)=4
    \end{align*}
\end{examplebox}